\begin{table}[htbp]
\centering
\small
\resizebox{\textwidth}{!}{%
\begin{tabular}{llllllllll}
\toprule
\textbf{path} & \textbf{scheme} & \textbf{judged against} & \textbf{mean confidence} & \textbf{accuracy} & \textbf{ECE [95\%]} & \textbf{MCE} & \textbf{Brier} & \textbf{constant Brier} & \textbf{Brier skill [95\%]} \\
\midrule
baseline & LOUO & self-declared skill (N/I/E) & 0.755 & 0.379 & 0.377 [0.237, 0.541] & 0.53 & 0.875 & 0.647 & -0.351 [-0.688, -0.092] \\
baseline & LOSO & self-declared skill (N/I/E) & 0.837 & 0.942 & 0.105 [0.057, 0.154] & 0.37 & 0.143 & 0.647 & 0.778 [0.698, 0.856] \\
jev & both & level implied by GRS & 0.696 & 0.301 & 0.395 [0.317, 0.494] & 0.55 & 1.021 & 0.713 & -0.433 [-0.588, -0.242] \\
\bottomrule
\end{tabular}
}
\caption{Calibration. Each path is judged against the target it was given, so the rows are not a head-to-head comparison. The constant forecast assigns every trial the observed class frequencies; a negative Brier skill score means the path is worse than that constant. Five equal-width bins, 103 pooled trials. Jev is never trained, so it has a single row. Per-bin counts are logged as \texttt{addendum\_reliability\_bin}; mean confidence is taken from \texttt{addendum\_mean\_confidence}, which postdates the calibration run.}
\label{tab:table4}
\end{table}
